%%%%%%%%%%%%%%%%%%%%%%%%%%%%%%%%%%%%%%%%%%%%%%%%%%%%%%%%%%%%%%%%%%%
\section{Limitations and Conclusions}
\label{sec:Conclusions}
%%%%%%%%%%%%%%%%%%%%%%%%%%%%%%%%%%%%%%%%%%%%%%%%%%%%%%%%%%%%%%%%%%%

\subsection{Limitations}
\label{subsec:Limitations}
Several qualifications apply to everything above.

\textbf{What we have not built:} Nothing here competes with a deployed
classifier, and no end-to-end jailbreak-blocking experiment was run. The
measured claims concern the fidelity and cost of stored-signature
matching, which is one component of the always-on stage described in
Sec.~\ref{sec:Introduction}.

\textbf{The label:} The experiments define deception as an output that
contradicts the known correct output. Deception in the stronger sense --
a model representing a claim as false while asserting it -- is harder to
label, and the available datasets control for it imperfectly.

\textbf{Granularity:} The grid results classify one activation per
response rather than per token, while the deployment target the hardware
serves is per-token and always-on. The watchlist reproduction of
Sec.~\ref{subsec:Watchlist} is per-token, which is part of why it is
reported here.

\textbf{Layer choice and model coverage:} Layer 33 is a tuned
hyperparameter of one model family at one scale, and re-tuning per model
is assumed. All grid results are one model. Re-tuning is also what makes
the bank a rewritten structure, which is why endurance appears in
Table~\ref{tab:devices}.

\textbf{Hardware numbers:} The comparison in Sec.~\ref{sec:Hardware} is
capacity and bandwidth arithmetic on vendor figures, with no contention
term for a monitor co-resident with live inference. The energy and
latency entries of Table~\ref{tab:costmodel} are first-order estimates
from literature values, not our measurements;
Sec.~\ref{subsec:NextMeasurements} lists what would replace them.

\textbf{Three metrics:} \auroc values in Sec.~\ref{sec:Results} are
cross-domain on the ten-domain grid; the first pipeline's numbers in
Sec.~\ref{subsec:FirstPipeline} are in-distribution on a four-domain
grid under a different protocol; and the readout numbers in
Sec.~\ref{sec:LargeK} are F1 and top-1 agreement against a
full-precision reference rather than against labels. None of the three is
comparable with the others cell by cell, and the paper does not place
them side by side.

\textbf{Digitized values:} the readout fidelity in
Fig.~\ref{fig:jspace} was digitized from result charts rather than read
from result files, and the library-size sweep carries three stated
implementation limitations (Sec.~\ref{subsec:JspaceScaling}). Both
should be regenerated from source before submission.

\subsection{Conclusions}
\label{subsec:Conclusions}
Always-on behavioral monitoring of a language model is a library-search
problem, and the variable that decides both its coverage and its cost is
$K$, the number of stored signatures compared against each activation. At
$K=1$, which is what is deployed today, the operation is nearly free and
no special hardware is warranted. Published work gives reasons to expect
$K$ to grow -- non-transfer across behaviors, steering wanting a
direction per behavior per domain per layer, feature dictionaries, and
whole-vocabulary readout, which fixes $K\!\approx\!10^{5}$ per layer by
construction -- and also gives reasons to doubt it, which is why the
question is worth measuring rather than asserting.

One methodological result comes before the fidelity numbers. Accuracy
alone does not establish that a memory-based detector is using its
memory: across five variants of an earlier pipeline, 66--100\% of the
accuracy of the two best-scoring configurations survived randomizing
every stored row, so the trained front end was doing most of the
classifying and the array was close to decoration. Reporting front-end
share alongside \auroc separates the two, and it is what moved this work
onto frozen projections.

What a \cam-resident monitor gives up can then be measured, and we
measure it on the two workloads at opposite ends of the range. For
deception detection at $10^{2}$--$10^{3}$ labeled rows, a 384-bit code
inside a plain array holds in-domain \auroc within 0.009 of
full-precision cosine similarity on the raw 8{,}192 numbers, at 0.920
(\signrp ternary) and 0.919 (thermometer) against 0.928; cross-domain,
the same codes give up 0.065 and 0.127. At a matched 384 bits the
ternary code is ahead of the thermometer code by 0.061 on average, which
is the opposite of what a preference for the cheaper cell would predict
and is the experiment we would run next at matched cell count rather
than matched code length. A two-stage shortlist recovers most of the
cross-domain gap in aggregate, but its fixed $k$ does not travel across
bank sizes, which argues for a threshold criterion over a top-$k$ one in
any array that has to serve banks of different sizes.

For whole-vocabulary readout at $10^{5}$ rows per layer by construction,
where the exact readout is its own ground truth and no labels are needed,
a 3-bit Euclidean code at $\sim$$10^{3}$ storage bits per row -- a
$128\times$ compression of an fp16 row -- reproduces 0.96 of the exact
top-1 selections and 0.81 of the per-row firing decisions. Reporting the
closest row survives compression better than reporting which rows cross a
threshold, by 0.15--0.2 at every code width, which is a matchline-mode
result rather than a cell-type one. Swept over library size, the bank
stays alive to $10^{4}$ rows while per-row fidelity falls from 0.755 to
0.43 -- and falls faster in full precision than at 3 bits, so what gets
harder with $K$ is the discrimination rather than the arithmetic.

On the cost side, monitoring traffic on a GPU crosses an H100's entire
HBM bandwidth at $K\!\approx\!204$K and a Jetson-class part's at
$K\!\approx\!17$K, while the same library at the measured code widths is
6--\qty{33}{\mega\byte} per layer and fits on a die. A first-order
estimate from literature figures puts the energy gap of the search
itself at five to six orders of magnitude -- roughly 67~mJ per token
per layer to stream the whole-vocabulary bank from HBM against
0.03--0.3~$\mu$J to compare it in place -- of which a factor of 64 is
the measured compression and the rest is the substrate
(Sec.~\ref{subsec:CostModelGap}). A library that is inexpensive to
hold and expensive to move is the condition under which searching it
in place is worth the silicon. The remaining work is short and
specific: Eva-CAM characterization at a stated node and a tuned GPU
baseline to turn the estimates into measurements
(Sec.~\ref{subsec:NextMeasurements}), and -- upstream of any hardware
-- establishing the minimum $K$ that adequate behavioral coverage
actually requires.

\begin{pnote}
Things I would add before submission, in order: {\bf (1)} the readout
fidelity numbers regenerated from result files rather than digitized, and
the library-size sweep re-run with the three fixes; {\bf (2)} the
matched-cell-count ternary vs.\ thermometer experiment;
{\bf (3)} the \textsc{est} entries of Table~\ref{tab:costmodel}
replaced with Eva-CAM output and a measured GPU baseline; {\bf (4)} the full
384 vs.\ 2{,}048-bit grids as figure panels; {\bf (5)} a second model,
even a small one, so the deception grid is not single-model. Item 1 is
the one a reviewer could actually catch us on.
\end{pnote}
